% PAGES: 226

% CONTINUES: closes the \begin{proof} of Theorem 7.5 part (i) left open at the end
% of sec07_c2_4.tex (PDF page 225, folio 209), mid-sentence right after "$X_1$,".
$X_2, \ldots, X_n$. Note that the covariance matrix $\Sigma_{\bfZ}$ of $Z_1$,
$Z_2$, \ldots, $Z_n$ is the identity matrix, i.e., $\Sigma_{\bfZ} = I$. Then, from
the Linear Transformation Property (Theorem 7.3) and using (7.91), we find that
\[
\Sigma_{\bfX} \;=\; M\Sigma_{\bfZ}M^t \;=\; MM^t \;=\; A.
\]
In other words, we found random variables $X_1$, $X_2$, \ldots, $X_n$ with
covariance matrix equal to the matrix $A$. We conclude that any symmetric
positive semidefinite matrix is a covariance matrix.

(ii) Recall from Lemma 7.4 that any correlation matrix is symmetric positive
semidefinite with main diagonal entries equal to $1$; see (7.11).

To prove that any symmetric positive semidefinite matrix with main diagonal
entries equal to $1$ is a correlation matrix, let $A$ be an $n \times n$
symmetric positive semidefinite with $A(i,i) = 1$ for all $i = 1:n$. We will find
random variables $X_1$, $X_2$, \ldots, $X_n$ with correlation matrix $\Omega_{\bfX}
= A$.

We showed above that there exist random variables $X_1$, $X_2$, \ldots, $X_n$
with covariance matrix $\Sigma_{\bfX} = A$. Since $A(i,i) = 1$ for all $i = 1:n$,
it follows that $\Sigma_{\bfX}(i,i) = 1$ for all $i = 1:n$, and, from Lemma 7.2,
we conclude that $\Omega_{\bfX} = \Sigma_{\bfX} = A$, which is what we wanted to
show.
\end{proof}

The method for finding normal random variables with a given covariance matrix
described above requires finding the eigenvalues of the given symmetric positive
semidefinite matrix. Note that this method is not used in practice if the given
matrix is symmetric positive definite, in which case the Cholesky decomposition
of the given covariance matrix is used; see section 7.5 for details.

\begin{examplex}
The $3 \times 3$ matrix
\begin{equation}
A \;=\; \begin{pmatrix} 1 & a & b \\ a & 1 & c \\ b & c & 1 \end{pmatrix}
\end{equation}
is a correlation matrix if and only if
\begin{equation}
-1 \le a,b,c \le 1 \quad \text{and} \quad \det(A) = 1+2abc-a^2-b^2-c^2 \ge 0.
\end{equation}

\par\medskip\noindent\textit{Solution:}\ Since all the main diagonal entries of
$A$ are equal to $1$, it follows from Theorem 7.5 that the matrix $A$ is a
correlation matrix if and only if it is symmetric positive semidefinite, which,
according to (5.64) is equivalent to
\begin{equation}
-1 \le a,b,c \le 1 \quad \text{and} \quad \det(A) = 1+2abc-a^2-b^2-c^2 \ge 0.
\end{equation}
\end{examplex}

\begin{examplex}
Find all the values of $\rho$ such that the matrix
\[
\Omega \;=\; \begin{pmatrix} 1 & 0.8 & 0.3 \\ 0.8 & 1 & \rho \\ 0.3 & \rho & 1 \end{pmatrix}
\]
is a correlation matrix.
% OPEN AT END OF PDF page 226 (folio 210): this \begin{examplex} is left
% deliberately UNCLOSED. The source's "Solution:" and the closing square do not
% appear until PDF page 227, outside this agent's assigned range (218-226).
% check_env_balance.py CANNOT see this gap: examplex is a manual run-in block
% (\textit{Example:} ... ends with $\square$ via \end{examplex}), not an
% environment the balance checker tracks structurally beyond its own
% begin/end pairing. THE NEXT AGENT (covering PDF 227+) MUST supply the missing
% "Solution:" text and close this block with \end{examplex} -- do not open a new
% \begin{examplex} for it.
